\documentclass[polish]{../uw-thesis}

\titleEN{Title of the doctoral dissertation}
\titlePL{Polski tytuł rozprawy doktorskiej}

\author{Imię Nazwisko}
%\place{Warszawa}
\date{Czerwiec 2026}

\faculty{Szkoła Doktorska Nauk Ścisłych i Przyrodniczych}
%\faculty{Międzydziedzinowa Szkoła Doktorska}
%\faculty{Wydział Matematyki, Informatyki i Mechaniki}

%\fieldPL{nauk ścisłych i przyrodniczych}

\disciplinePL{informatyka}
%\disciplinePL{matematyka}

\addSupervisor{Prof.~Główny Promotor}
\addSupervisor{Prof.~Jego Kolega}
\addCosupervisor{Dr~Pomocniczy Promotor}
\addCosupervisor{Dr~Jego Ziomeczek}
\addCosupervisor{Mgr~Ich Koleś}

\keywords{słowo kluczowe 1; słowo kluczowe 2; słowo kluczowe 3}

\abstractPL{Tutaj wstaw streszczenie po polsku.}
\abstractEN{Put the English abstract here.}

\newtheorem{theorem}{Twierdzenie}[chapter]
\newtheorem{lemma}[theorem]{Lemat}

\begin{document}
\makefrontmatter

\chapter{Wprowadzenie}
Tutaj rozpoczyna się rozprawa.

\section{Sekcja pierwsza}

\begin{theorem}
    Proste twierdzenie. Możesz używać środowiska \texttt{theorem} dla twierdzeń, lem, propozycji itp.
\end{theorem}

\cite{tarski_decision}

\section{Sekcja druga}

\begin{theorem}
    To jest drugie twierdzenie!
\end{theorem}

\begin{lemma}
To jest próba lematu.
\end{lemma}

\chapter{Zawartość}

To jest drugi rozdział.

\begin{lemma}
To jest próba lematu.
\end{lemma}

\begin{lemma}
To jest próba lematu.
\end{lemma}

\bibliographystyle{alpha}
\bibliography{test.bib}
\end{document}
